% ---------------------------------------------------------------------------
% SECTION 7: SOFTWARE
% ---------------------------------------------------------------------------
% ---------------------------------------------------------------------------
% SECTION 7: SOFTWARE (one page; details in Appendix B)
% ---------------------------------------------------------------------------
\section{Software Implementation}\label{sec:software}

The methods are implemented in the open-source R package
\texttt{adaptsimex}, whose main entry point is

\begin{center}
\small
\begin{minipage}{0.9\textwidth}
\texttt{adaptive\_simex(formula, data, me\_vars, sigma\_error,\\
\quad se\_method = c("sandwich", "bootstrap"), seed = NULL)}
\end{minipage}
\end{center}

\noindent where \texttt{me\_vars} names the error-contaminated
columns and \texttt{sigma\_error} their standard deviations. The
function cross-validates $(\lambda_{\max},\text{degree})$ with the
one-standard-error rule, fits the two-stage extrapolated estimator,
and returns coefficients with sandwich or bootstrap standard errors,
the CV score table, and the fitted SIMEX curve for the diagnostic
plots. The analytic sandwich of Section~\ref{sec:var} is the default;
the pairs-bootstrap is recommended at small samples
(Section~\ref{sec:sims-val}). The package ships with unit tests, a
vignette reproducing the attenuation-correction comparison on
simulated data, and S3 methods (\texttt{summary}, \texttt{confint},
\texttt{plot}). Every number in Sections~\ref{sec:sims}
and~\ref{sec:app} was produced by the same code path (a documented
Python harness that ships with the replication archive and matches
the R implementation to numerical precision), so the archive
regenerates the paper exactly. Package structure, the full call
interface, and the computational tricks that bring a full A-SIMEX fit
down to seconds are documented in the Supplementary Material (Section~S2);
timing evidence appears in the Supplementary Material (Section~S2).

% ---------------------------------------------------------------------------
\section{Discussion}\label{sec:discussion}

\subsection{Summary of Contributions}

This paper makes SIMEX for Poisson regression automatic, gives it
complete and corrected asymptotic theory, and reports its
finite-sample behavior at replication standard. Three findings that we
believe are of independent interest emerged from holding the method to
this standard: (i)~held-out prediction on the observed scale cannot
select the extrapolation that bias correction requires, because its
population minimizer is the naive pseudo-parameter; (ii)~the variance
of the extrapolated functional is governed by the norm of the
extrapolation weights ($\sum_j|c_j|\approx7.8$ for the standard
quadratic extrapolant), which quantifies a substantial and previously
unappreciated efficiency premium; and (iii)~the analytic
data-plus-simulation sandwich makes inference nearly free, and its
$66$--$91\%$ recovery of the bootstrap variance delineates what part
of the finite-sample noncoverage is variance miscalibration and what
part is residual extrapolation bias.

\subsection{Comparison with Alternatives}

\emph{Versus standard SIMEX.} A-SIMEX removes the ad-hoc choices and
adds the variance machinery; in the simulations the two track each
other closely, with A-SIMEX (1-SE rule) marginally better at small
error and marginally worse at large error, where the parsimony rule
reinforces the CV's under-extrapolation. The honest summary is that
the tuning layer is not where the accuracy is won or lost; the
extrapolation bias is.

\emph{Versus the corrected score \citep{nakamura1990}.} The corrected
score is consistent for the Gaussian-error pseudo-score without
requiring a model for $\Pr(X\mid W)$: it requires only the error
distribution, as SIMEX does, and it is dramatically cheaper
($<0.05$\,s at $p=50$). Its weaknesses, quantified in
Section~\ref{sec:sims}, are numerical: Newton divergence in about
$1\%$ of fits under misspecification, RMSE exploding with $p$ when
many coefficients are weakly identified ($0.13$ at $p=5$ to $2.72$ at
$p=50$), and instability at small $n$ with large error. SIMEX-type
methods are more stable in exactly these regimes, at a computational
premium. The two are complements: the corrected score when $p$ is
large and the error small, SIMEX when stability matters.

\emph{Versus deconvolution and functional methods
\citep{fan1993,delaigle2008}.} These remain computationally intensive
and ill-posed; A-SIMEX scales to $p=50$ in seconds. A fair comparison
at their level of generality would be valuable but is beyond our
scope.

\subsection{Limitations and Caveats}\label{sec:caveats}

\begin{enumerate}\itemsep3pt
\item \textbf{Error variance known or estimated.} The theory
      conditions on $\bSig_\epsilon$. The replicate-based estimator
      performs indistinguishably in Section~\ref{sec:sims-sigma2};
      the estimated-variance asymptotics (uncertainty propagation
      through the extrapolation) are not developed here.
\item \textbf{Classical error.} The NHANES application makes the
      assumption auditable: the corrected estimates lie $1.5$ standard
      errors \emph{above} the measured-weight anchor, the signature of
      nonclassical (truth-correlated) self-report error, which no
      classical-error method removes. Users with validation data
      should always fit the anchor when it exists.
\item \textbf{Residual extrapolation bias at large error.} At
      $\sigma^2_\epsilon/\Var(X)\gtrsim0.5$, polynomial SIMEX of any
      flavor retains bias of the same order as the naive attenuation
      it removes; coverage of nominal $95\%$ intervals degrades
      accordingly. Remedies (deconvolution-aware selection;
      variance-penalized degree selection) are stated in
      Section~\ref{sec:sims-val} but not developed.
\item \textbf{First-order variance estimator.} The sandwich
      under-states the finite-sample variance by $10$--$34\%$ at
      $n\le100$; bootstrap standard errors are the recommended
      small-sample alternative.
\item \textbf{High dimensions.} Repeated GLM fitting remains
      $O(np^2)$ per pseudo-sample; regularized SIMEX would be needed
      for $p$ in the hundreds.
\end{enumerate}

\subsection{Extensions and Future Work}\label{sec:extensions}

Beyond Berkson error, imputation with validation data, alternative
GLMs, time-series and spatial dependence, and regularized extensions
(all still open), the present results add two concrete problems:
(i)~\emph{deconvolution-aware model selection}: constructing a
held-out criterion whose population minimizer is
$\beta(-1)=\bet_0$ rather than $\beta(1)$, for example by evaluating
the held-out deviance at pseudo-deconvolved covariates simulated under
each candidate hyperparameter; our CV-bias finding provides both the
motivation and the diagnostic; and (ii)~\emph{variance-aware degree
selection}: enlarging the degree grid while penalizing the
extrapolation variance $\dot{\bm C}_\beta$ of Theorem~\ref{thm:are},
in the spirit of model-selection penalties, which would automate the
bias--variance trade-off that practitioners currently negotiate by
hand.

% ---------------------------------------------------------------------------
% SECTION 9: CONCLUSION (revised)
% ---------------------------------------------------------------------------
\section{Conclusion}\label{sec:conclusion}

Measurement error in covariates is pervasive, and the naive Poisson
fit is inefficient and misleading: its attenuation grows
linearly in the error variance and its intervals collapse, as our
simulations document cell by cell. SIMEX remains the most practical
correction framework for this problem, and this paper makes it
automatic, gives it correct and complete asymptotics, equips it with a
closed-form variance estimator, and reports, at replication
standard with all seeds fixed, exactly where it succeeds (small to
moderate error, with valid inference and RMSE ties with the best
competitor), what it costs (a quantified efficiency premium relative
to the oracle), and where it fails (large error, where residual
extrapolation bias dominates and observed-scale cross-validation
under-extrapolates by construction). The NHANES application shows the
workflow we recommend in practice: estimate the error variance from
measurable discrepancies, correct, validate against any available
anchor, and report the sensitivity profile over the error variance.
The open-source implementation and the fully reproducible archive are
offered so that both the method and the standard of evidence can be
checked and extended.

% ---------------------------------------------------------------------------
% REFERENCES (APA)
% ---------------------------------------------------------------------------
\section*{Data and code availability}\label{sec:data}

The NHANES public-use files used in Sections~\ref{sec:app}
and~\ref{sec:transfer} (2017--2018: DEMO\_J, BMX\_J, WHQ\_J,
RXQ\_RX\_J; 2015--2016: the corresponding H files) are distributed by
the National Center for Health Statistics and are downloaded by the
replication archive. The archive contains the R
package, the simulation and application code, all raw Monte Carlo
output, and a master seed (20260905) from which every table and figure
in the paper regenerates exactly; see the README.

\begin{thebibliography}{25}
\providecommand{\natexlab}[1]{#1}

\bibitem[Apanasovich et al.(2009)]{apanasovich2009}
Apanasovich, T.~V., Carroll, R.~J., \& Maity, A. (2009). SIMEX and
standard error estimation in semiparametric measurement error models.
\emph{Electronic Journal of Statistics}, \emph{3}, 558--585. % page range to verify

\bibitem[Berkson(1950)]{berkson1950}
Berkson, J. (1950). Are there two regressions? \emph{Journal of the
American Statistical Association}, \emph{45}(250), 164--180.

\bibitem[Berry et al.(2002)]{berry2002}
Berry, S.~M., Carroll, R.~J., \& Ruppert, D. (2002). Bayesian smoothing
and regression splines for measurement error problems. \emph{Journal
of the American Statistical Association}, \emph{97}(457), 160--169.

\bibitem[Bound et al.(2001)]{bound2001}
Bound, J., Brown, C., \& Mathiowetz, N. (2001). Measurement error in
survey data. In J.~J. Heckman \& E.~Leamer (Eds.), \emph{Handbook of
Econometrics} (Vol.~5, pp.~3705--3843). Elsevier.

\bibitem[Cameron and Trivedi(2013)]{cameron2013}
Cameron, A.~C., \& Trivedi, P.~K. (2013). \emph{Regression analysis of
count data} (2nd~ed.). Cambridge University Press.

\bibitem[Carroll et al.(2006)]{carroll2006}
Carroll, R.~J., Ruppert, D., Stefanski, L.~A., \& Crainiceanu, C.~M.
(2006). \emph{Measurement error in nonlinear models: A modern
perspective} (2nd~ed.). Chapman \& Hall/CRC.

\bibitem[Chen et al.(2011)]{chen2011}
Chen, X., Hong, H., \& Nekipelov, D. (2011). Nonlinear models with
measurement error: A brief survey. \emph{Journal of Economic
Literature}, \emph{49}(4), 906--937.

\bibitem[Cook and Stefanski(1994)]{cook1994}
Cook, J.~R., \& Stefanski, L.~A. (1994). Simulation-extrapolation
estimation in parametric measurement error models. \emph{Journal of
the American Statistical Association}, \emph{89}(428), 1314--1328.

\bibitem[Delaigle and Meister(2008)]{delaigle2008}
Delaigle, A., \& Meister, A. (2008). Density estimation from
heteroscedastic data. \emph{Journal of the American Statistical
Association}, \emph{103}(482), 726--736.

\bibitem[Fan(1991)]{fan1991}
Fan, J. (1991). On the optimal rates of convergence for nonparametric
deconvolution problems. \emph{The Annals of Statistics}, \emph{19}(3),
1257--1272.

\bibitem[Fan and Truong(1993)]{fan1993}
Fan, J., \& Truong, Y.~K. (1993). Nonparametric regression with errors
in variables. \emph{The Annals of Statistics}, \emph{21}(3),
1900--1925.

\bibitem[Fuller(1987)]{fuller1987}
Fuller, W.~A. (1987). \emph{Measurement error models}. John Wiley \&
Sons.

\bibitem[Griliches(1986)]{griliches1986}
Griliches, Z. (1986). Economic data issues. In Z.~Griliches \&
M.~D. Intriligator (Eds.), \emph{Handbook of Econometrics} (Vol.~3,
pp.~1465--1514). North-Holland. % page range to verify

\bibitem[Hausman(2001)]{hausman2001}
Hausman, J.~A. (2001). Mismeasured variables in econometric analysis:
Problems and solutions. \emph{Journal of Econometrics}, \emph{101},
1--39. % page range to verify

\bibitem[Hu(2008)]{hu2008}
Hu, Y. (2008). Identification and estimation of nonlinear models with
misclassification error using instrumental variables: A general
solution. \emph{Econometrica}, \emph{76}(1), 1--33. % page range to verify

\bibitem[K\"uchenhoff et al.(2006)]{kuechenhoff2006}
K\"uchenhoff, H., Mwalili, S.~M., \& Lesaffre, E. (2006). A general
method for dealing with misclassification in regression: The
misclassification SIMEX. \emph{Biometrics}, \emph{62}(1), 85--94.

\bibitem[Lewbel(1997)]{lewbel1997}
Lewbel, A. (1997). Constructing instruments for regressions with
measurement error when no additional data are available, with an
application to patents and R\&D. \emph{Econometrica}, \emph{65}(5),
1201--1213.

\bibitem[Mahajan(2006)]{mahajan2006}
Mahajan, A. (2006). Identification and estimation of regression models
with misclassification. \emph{Econometrica}, \emph{74}(3), 645--678.

\bibitem[Mullahy(1986)]{mullahy1986}
Mullahy, J. (1986). Specification and testing of some modified count
data models. \emph{Journal of Econometrics}, \emph{33}(3), 341--365.

\bibitem[Nakamura(1990)]{nakamura1990}
Nakamura, T. (1990). Corrected score function for errors-in-variables
models: Methodology and application to generalized linear models.
\emph{Biometrika}, \emph{77}(1), 127--137.

\bibitem[National Center for Health Statistics(2020)]{nhanes-rxq}
National Center for Health Statistics. (2020). \emph{NHANES 2017--2018:
Prescription medications file documentation} (RXQ\_RX\_J). Centers for
Disease Control and Prevention, Hyattsville, MD.

\bibitem[Prentice(1982)]{prentice1982}
Prentice, R.~L. (1982). Covariate measurement errors and parameter
estimation in a failure time regression model. \emph{Biometrics},
\emph{38}(3), 601--612.

\bibitem[Qin(2000)]{qin2000}
Qin, J. (2000). Combining parametric and empirical likelihoods.
\emph{Biometrika}, \emph{87}(3), 484--490.

\bibitem[Richardson and Gilks(1993)]{richardson1993}
Richardson, S., \& Gilks, W.~R. (1993). Conditional independence models
for epidemiological studies with covariate measurement error.
\emph{Biometrics}, \emph{49}(3), 901--911. % page range to verify

\bibitem[Rivlin(1969)]{rivlin1969}
Rivlin, T.~J. (1969). \emph{An introduction to the approximation of
functions}. Blaisdell Publishing Company.

\bibitem[Rosner et al.(1989)]{rosner1989}
Rosner, B., Willett, W.~C., \& Spiegelman, D. (1989). Correction of
logistic regression relative risk estimates and confidence intervals
for systematic underreporting of intake. \emph{Statistics in Medicine},
\emph{8}(9), 1051--1069.

\bibitem[Schennach(2004)]{schennach2004}
Schennach, S.~M. (2004). Estimation of continuous models with errors in
variables. \emph{Econometrica}, \emph{72}(1), 337--341. % page range to verify

\bibitem[Stefanski and Cook(1995)]{stefanski1995}
Stefanski, L.~A., \& Cook, J.~R. (1995). Simulation-extrapolation: The
measurement error jackknife. \emph{Journal of the American Statistical
Association}, \emph{90}(432), 1247--1256.

\bibitem[Trefethen(2013)]{trefethen2013}
Trefethen, L.~N. (2013). \emph{Approximation theory and approximation
practice}. Society for Industrial and Applied Mathematics.

\bibitem[van der Vaart(1998)]{vdv1998}
van der Vaart, A.~W. (1998). \emph{Asymptotic statistics}. Cambridge
University Press.

\bibitem[van der Vaart and Wellner(1996)]{vdv1996}
van der Vaart, A.~W., \& Wellner, J.~A. (1996). \emph{Weak convergence
and empirical processes: With applications to statistics}. Springer.

\bibitem[Winkelmann(2008)]{winkelmann2008}
Winkelmann, R. (2008). \emph{Econometric analysis of count data}
(5th~ed.). Springer.
\end{thebibliography}
